\section{Actual production and Houghton coercivity}
\label{app:houghton-coercivity}

\subsection{The root-energy dictionary}

\begin{proposition}[Actual-production dictionary]\label{prop:actual-dirichlet}
For the literal smoothed slots $T\rho=\sum_sw_sV_s\rho V_s^*$,
with $m\leq174$, $V_0=I$ and $w_0\geq1/d$, put
$\mathfrak D=\sum_sw_s\|V_s\sqrt\rho V_s^*-\sqrt\rho\|_F^2$.
For every source, with $a=S(T\rho)-S(\rho)$,
\[
 \mathfrak D\leq(d+1)\ln2\,a,\qquad
 a\leq h_2(t)+t\log_2(m-1),\quad
 t=\min\{\mathfrak D,(m-1)/m\}.
\]
On the same feasible physical rounds, optimizing $S+n\mathfrak D$
and $S+na$ gives equivalent costs up to $O(\log(en))$ and an
additive constant. Error and leakage constraints are unchanged.
\end{proposition}

\begin{proof}[Proof of Proposition~\ref{prop:actual-dirichlet}]
Put $\xi=\sqrt\rho$, $\xi_s=V_s\xi V_s^*$,
$\bar\xi=\sum_sw_s\xi_s$, and $\sigma=T\rho$.
The Holevo identity gives $a=\sum_sw_sD_2(\rho_s\Vert\sigma)$.
For states $P,Q$,
\[
 D_{\rm nat}(P\Vert Q)\geq\|\sqrt P-\sqrt Q\|_F^2.
\]
Indeed spectral overlap probabilities reduce the left side to a
classical relative entropy, and Jensen gives
$D_{\rm nat}\geq-2\ln\Tr\sqrt P\sqrt Q
\geq2-2\Tr\sqrt P\sqrt Q$; zero eigenvalues follow by limits.
The variance identity and minimization over a common centre give
\[
 v:=\sum_sw_s\|\xi_s-\bar\xi\|_F^2
     =1-\|\bar\xi\|_F^2\leq(\ln2)a.
\]
The identity slot implies $w_0\|\xi-\bar\xi\|_F^2\leq v$, hence
$\mathfrak D=v+\|\bar\xi-\xi\|_F^2\leq(d+1)(\ln2)a$.

For the converse, form the pure-state ensemble
$\Omega=\sum_sw_s|\operatorname{vec}\xi_s\rangle
\langle\operatorname{vec}\xi_s|$. Partial trace maps it to the
ensemble $\rho_s$, so Holevo data processing gives $a\leq S(\Omega)$.
Its rank is at most $m$ and
$\Omega\geq|\operatorname{vec}\bar\xi\rangle
\langle\operatorname{vec}\bar\xi|$, so
$\lambda_{\max}(\Omega)\geq1-v\geq1-\mathfrak D$.
Maximizing entropy with this eigenvalue constraint gives the stated
binary-entropy bound. Zero weights can be omitted; $m=1$ is trivial.

For completeness put $y=\log_2(en)$ and $C_0=(d+1)\ln2$.
The first inequality gives
$S+n\mathfrak D\leq\max(1,C_0)(S+na)$.
The second, using $h_2(t)\leq t\log_2(e/t)$, gives
$S+na\leq C_m y(1+S+n\mathfrak D)$ for a fixed $C_m$.
If $\mathfrak D\leq1/n$, set $x=n\mathfrak D\leq1$ and use
$x\log_2(en/x)\leq y+1/(e\ln2)$.
If $1/n<\mathfrak D<(m-1)/m$, use
$\log_2(e/\mathfrak D)\leq y$; above this threshold use
$a\leq\log_2m\leq2(\log_2m)\mathfrak D$.
Taking infima over the same feasible family proves the profile claim.
\end{proof}

\subsection{The exact leakage coefficient}

Here and below $\Pi_k$ denotes the smoothed model's bath
representation, $W=r_4$, and
$X=(\Pi_k(W)-I)^*(\Pi_k(W)-I)$.
Inspecting the canonical slots gives, for every source,
\begin{equation}\label{eq:actual-houghton-leakage}
 \ell=\frac{488}{225d}\Tr\rho X.
\end{equation}
The calculation is weighted least squares, not a defect inequality.
The identity class has total squared coefficient $45/2$:
one exceptional closing slot of weight $1/2$ versus weight $22$
with the usual endpoint. The $\alpha$ class has total $25/2$:
exceptional weight $2$ versus usual weight $21/2$.
The other exceptional word $\sfa\sfb\sfa^{-1}=r_4\alpha\sfb$
is alone in its target label class and contributes no residual.
The flux phases act on coordinates numbered at least five, so they
commute with $\alpha$. Thus the two variances multiply the same $X$:
\[
 \Lambda=\left(\frac{(1/2)22}{45/2}
                  +\frac{2(21/2)}{25/2}\right)X
          =\frac{488}{225}X .
\]
All other slot classes are consistent.
Let $Q_\gamma$ be the projection onto the nontrivial eigenvalues
of the order-three operator $\Pi_k(\gamma)$. The anchor gives
$|\Tr\rho\Pi_k(\gamma)|\leq2u\leq1/8$, so
\begin{equation}\label{eq:actual-houghton-charge}
 \Tr\rho Q_\gamma\geq7/12,
\end{equation}
since $\operatorname{Re}\Tr\rho\Pi_k(\gamma)
=1-\tfrac32\Tr\rho Q_\gamma$.

\subsection{A representation-uniform gap at fixed clock size}

\begin{theorem}[Fixed-clock coercivity]\label{thm:fixed-clock}
For each fixed $M\geq48$ in the literal rational constant-field
smoothed model there is $c_M>0$, independent of the flavour count $k$,
such that every source with $u\leq1/16$ satisfies $a+\ell\geq c_M$.
Consequently $\ell=O(1/n)$ forces $S+na=\Omega_M(n)$.
\end{theorem}

We give the compact-group argument for Theorem~\ref{thm:fixed-clock}.
Its only geometric inputs are the actual unitary model of
\cite[Section~6]{DouglasMghirbiA}, specialized to its rational
constant fields. At fixed $M\geq48$, put $L=4M+4$.
The one-particle generators lie in
\[
 G=U(L)^{M^2}\rtimes\mathbb Z_M^2
       \ \subset U(\C^{M^2}\otimes\C^L).
\]
Their links are permutations times $e^{i\phi Q_{\rm pair}}$.
The tree curl construction with fields $\pm\pi/|R|$ makes $\phi$
a rational multiple of $\pi$, so all their entries are algebraic.
The words $r_1,r_2,r_3,r_5$ are exact, $\gamma$ is a three-cycle,
and $W=r_4$ has disjoint-plane phases. In particular
\[
 W^q=I,\quad q=2\,\operatorname{lcm}_R|R|,\qquad
 \|W-I\|_{\rm op}\leq\pi/m_M<1/2,\quad
 m_M=1+(M-1)(M-2)/2.
\]
For each $k$, the model on
$\C^{M^2}\otimes\Lambda(\C^L)^{\otimes k}$
is a continuous unitary representation of this same $G$.
It is not the exterior power of the full one-particle clock space.

Let $\Gamma=\langle\sfa,\sfb,\alpha\rangle$, $K=\overline\Gamma$,
and let $\Gamma_N$ be the abstract normal closure of $W$ in $\Gamma$.
Then $N=\overline{\Gamma_N}$ is a closed normal subgroup of $K$.
We first prove that $N^0$ is semisimple.
The torus $Z=Z(N^0)^0$ is characteristic and hence $K$-invariant.
Every $K$-conjugate of $W$ acts on its Lie algebra with
\[
 \|\operatorname{Ad}(W)-I\|_{F\to F}
       \leq2\|W-I\|_{\rm op}<1.
\]
This is an order-$q$ torus-lattice automorphism. A nonidentity
finite-order integral automorphism cannot have every eigenvalue
within distance strictly less than one from $1$: a nontrivial
root of unity $\zeta$ brings all its Galois conjugates in the
integer characteristic polynomial, while the nonzero integer
$\prod_{\zeta'}(1-\zeta')$ would then have absolute value less than
one. Thus every conjugate centralizes $Z$, and density shows that
$N$ does too.

The compact abelian Lie quotient $N/\overline{[N,N]}$ has exponent
dividing $q$, since it is generated densely by these order-$q$
conjugates; it is therefore finite. Put $D=[N^0,N^0]$, a closed
semisimple subgroup. In $N/D$, the connected component is the
image of $Z$, hence is central and has finite index.
Schur's theorem gives a finite commutator subgroup
\cite{Yadav2016}. Its abelianization is also finite, as a quotient
of the preceding exponent-$q$ abelian group. Thus $N/D$ is finite,
and its connected component is trivial. This proves $N^0=D$.

Next $\gamma\in N$. The quotient $K/N$ satisfies the $H_3$
presentation. Every finite-dimensional unitary representation
of $H_3$ kills its subgroup $\operatorname{Alt}_{\rm fin}$:
simplicity would otherwise give an injection, whereas this
group contains elementary abelian $2$-groups of arbitrarily
large rank and commuting involutions in $U(p)$ have rank at most
$p$. Continuous finite-dimensional representations of a compact
group separate its points. They therefore kill the image of
$\gamma$ in $K/N$, proving the assertion.

The dense subgroup $\Gamma_N\cap N^0$ consists of algebraic
matrices; $N^0$ is open in $N$.
If $N^0$ is nontrivial, Breuillard--Gelander's theorem supplies
two elements of this subgroup generating a dense subgroup of
$N^0$ \cite[Theorem~1.1]{BreuillardGelander2003}.
Their adjoint matrices have algebraic entries in an algebraic
Lie-algebra basis. Here is the basis justification, to exclude
an implicit rationality assumption on a closed subgroup.
A compact matrix group is real Zariski closed: polynomial
approximation followed by averaging over the group separates
any exterior point from the group by an invariant polynomial.
At each bounded degree its vanishing polynomials are the kernel
of evaluations on the dense algebraic subgroup. Finite-rank
linear algebra selects finitely many algebraic evaluation rows,
so this kernel has an algebraic coefficient basis.
Noetherianity selects finitely many of these defining equations.
Taking their tangent equations at the identity yields an
algebraic basis of the Lie algebra; compact real algebraic groups
are smooth. Conjugation by the selected algebraic matrices is
algebraic in this basis.

Apply He--de Saxce's theorem to a symmetric lazy probability
measure on these generators and their inverses
\cite[Theorem~1.4]{HeDeSaxce2021}. It has a spectral gap on
$L^2_0(N^0)$. Peter--Weyl then gives the same gap in every
continuous unitary representation after removing $N^0$-invariants.
Add finitely many representatives in $\Gamma_N$ of $N/N^0$.
The projection onto $N^0$-invariants commutes with $N$; on its
range the finite quotient has a uniform gap.
Combining the two gaps gives a finite set $F_M\subset\Gamma_N$
and a representation-uniform constant. If $N^0$ is trivial,
the finite-group gap gives the same conclusion directly.
Every member of $F_M$ is a finite word in $\Gamma$-conjugates of
$W$ and $W^{-1}$. All choices and constants depend on $M$,
not on a representation or its dimension.

On Hilbert--Schmidt matrices use
$\zeta_k(x,y)A=\Pi_k(x)A\Pi_k(y)^*$.
Left multiplication by a conjugate of $W$ is
$(g,g)(W,e)(g^{-1},g^{-1})$ in this representation;
there is an analogous right formula.
Telescoping each fixed word bounds its squared displacement by
its length times the sum of the squared elementary displacements.
Diagonal $\sfa,\sfb,\alpha$ displacements are bounded by $d$ times
the actual slot Dirichlet energy, since their singleton weights
are at least $1/d$; inverse displacements have the same norm.
The remaining elementary energies are left and right $W-I$.
Apply the finite-set gap separately to the left and right
representations of $N$. Their invariant projections $P_L,P_R$
commute, and distance to $P_LP_R$ is at most the sum of the two
squared distances. Since $\gamma\in N$, its left and right
nontrivial spectral projections annihilate $P_LP_R$. We obtain
\begin{equation}\label{eq:fixed-clock-operator}
 \begin{aligned}
 \frac{\|Q_{\gamma,k}A\|_F^2+\|AQ_{\gamma,k}\|_F^2}{2}
 &\leq C_M\left[\sum_sw_s\|V_sAV_s^*-A\|_F^2\right.\\
 &\quad\left.+
 \frac{\|(\Pi_k(W)-I)A\|_F^2+\|A(\Pi_k(W)-I)\|_F^2}{2}
 \right].
 \end{aligned}
\end{equation}
Set $A=\sqrt\rho$. The two symmetrized terms become
$\Tr\rho Q_\gamma$ and $\Tr\rho X$, respectively.
Equations~\eqref{eq:actual-houghton-leakage}--\eqref{eq:actual-houghton-charge}
and Proposition~\ref{prop:actual-dirichlet} give
\[
 7/12\leq C_M\left[(d+1)(\ln2)a+
                      \frac{225d}{488}\ell\right],
\]
proving the fixed-clock theorem. No invariance or entropy condition
on $\rho$ has been introduced.

\subsection{Entropy without source symmetry}

\begin{theorem}[Entropy without source symmetry]\label{thm:houghton-entropy}
In the literal rational constant-field smoothed model, every $M\geq48$,
$k\geq1$ and arbitrary coherent source with $u\leq1/16$ satisfies,
with a universal $c>0$,
\[
 S(\rho)+1\geq c\min\{M,a(\rho)^{-1/2}\},\qquad
 S(\rho)+na(\rho)+1\geq c\min\{M,n^{1/3}\}\quad(n\geq1).
\]
At $a=0$ the first minimum is $M$.
\end{theorem}

\begin{proof}[Proof of Theorem~\ref{thm:houghton-entropy}]
Write $\pi$ for the round's usual free-group representation and
put $c_g=2(104\ln2)^{1/2}$. Lemma~\ref{lem:centrality} and
word telescoping give
$\chi(\pi(w))\leq |w|c_g\sqrt a$ and
$\tfrac12\|\rho-\pi(w)\rho\pi(w)^*\|_1\leq\chi(\pi(w))$.
The same Holevo argument after tracing to either clock coordinate
bounds its cyclic root energy $E_i$ by $c_g^2a$.
Cauchy--Schwarz gives $\max p_i\leq1/M+2\sqrt{E_i}$.
Hence $G_R=\{0\leq i,j\leq M-R-1\}$ has
$p(G_R^c)\leq2R/M+4Rc_g\sqrt a$.

We first verify the exact structure on $G_R$, for $5\leq R<M/4$.
At $c=(i_0,j_0)$, a forward horizontal prefix of length at most
$R-1$ has edge starts at most $M-3$ and endpoints at most $M-2$.
It meets neither the seam nor the preliminary exact horizontal move
of~\cite[Section~6.2]{DouglasMghirbiA}. Its base links are
$c_{2M+2-i-j_0}$. A horizontal edge is interior to a region only
when both incident plaquettes belong to it. For $j_0=0$ none of
these edges carries a rotation; for $j_0\geq1$ the bulk condition is
$i+j_0\leq t+M-8$, and the $T$ condition is $i+j_0\leq M-3$.
Thus a region eligible later was eligible earlier. Actual potentials
may be sparse; only eligibility is monotone.
Pulling each phase plane back through the base prefix gives its
fixed starting pair $(x,x+M+i_0)$, where
$x=t+M-1-i_0-j_0$ for $5\leq t\leq M+3$, and
$x=M+3-i_0-j_0$ for $T$. These pairs are disjoint.

In column-operator notation the prefix factors as $P_j=B_jR_j$,
with $R_j$ a product of commuting rotations on these pairs.
The usual representation takes adjoints of the right-action
operators, so the one-particle holonomy $J_j=P_j^*\tau P_j$ of
$\pi(\sfa^j\alpha\sfa^{-j})$ has normal
$R_j^*(e_{j+1}-e_{j+2})/\sqrt2$.
The far partners exceed $R+1$. Nonadjacent normals therefore have
disjoint support; adjacent overlaps have modulus at most $1/2$.
The exact conjugates of $r_2$ force that modulus to be $1/2$.
Phase adjustment gives the standard $A_R$ root Gram matrix, proving
an exact $S_{R+1}$ representation. Quantization preserves it for
every flavour count.

Partition this chain into $t_R=\lfloor(R+1)/3\rfloor$ commuting
$S_3$ blocks. Their three-cycles are
$\gamma_b=\sfa^{3b}\gamma\sfa^{-3b}$.
The original anchor gives $\Tr\rho Q_\gamma\geq7/12$.
Transport and restriction to $G_R$ leave each block weighted charge
at least $7/12-2R/M-5Rc_g\sqrt a$.
Six representatives for each block have word lengths at most
$2R+5\leq3R$. Their canonical root pure-state ensemble has rank
at most six and an eigenvalue at least $1-x$,
$x=9R^2c_g^2a$. Partial trace bounds the mixed-state Holevo gain by
$f_6(x)=h_2(x)+x\log_2 5$ for $x\leq5/6$.

Clock pinching and Holevo data processing bound the weighted
conditional gains on $G_R$ by $f_6(x)$. On these clocks only, each
previous block twirl commutes with every conjugation in a later
block, so it cannot increase that block's gain. The final state is
flat on the irreducible factors of $S_3^{t_R}$, with arbitrary
multiplicities. Only its two-dimensional standard factors have
nontrivial three-cycle charge; those charge projections are central
and preserved. Final conditional entropy is therefore at least the
sum of charges. Coherent clock copying and Araki--Lieb give
$S(\rho)\geq\sum_cp_cS(\rho_c)$, hence
\[
 S(\rho)\geq t_R\left[\frac7{12}-\frac{2R}{M}
       -5Rc_g\sqrt a-f_6(9R^2c_g^2a)\right].
\]
No global commutation on the excluded clocks is required.

Put $K=\min\{M,1/(c_g\sqrt a)\}$ and $R=\lfloor K/128\rfloor$.
For $R\geq5$, the bracket exceeds $7/16$, since
$7/12-7/128-f_6(9/16384)>0.5206$; also $t_R\geq R/4$.
Thus $S>R/10$ and $S+1\geq K/1280$.
For $R<5$, $K<640$ and $S+1\geq1$ suffices.
This proves the first inequality with $c=1/(1280c_g)$.
For the second, $a\geq n^{-2/3}$ gives $na\geq n^{1/3}$;
otherwise the first inequality applies with $a^{-1/2}>n^{1/3}$.
\end{proof}

\subsection{A fixed local group and the fourth-power leakage bound}

\begin{theorem}[Leakage without source symmetry]\label{thm:houghton-flux}
For the literal rational constant-field smoothed round, all
$M\geq48$, $k\geq1$ and arbitrary coherent sources with $u\leq1/16$
satisfy, with constants $c_0,c>0$ independent of $M,k,\rho$,
\[
 M^2a+M^4\ell\geq c_0,\qquad
 S+na+1\geq c\min\{n^{1/3},\ell^{-1/4}\}\quad(n\geq1).
\]
At zero leakage the second argument is infinite.
\end{theorem}

\begin{proof}[Proof of Theorem~\ref{thm:houghton-flux}]
Put $d=104$, $c_g=2\sqrt{104\ln2}$ and
$\sigma=\Delta_{\rm clock}\rho=\bigoplus_cp_c\rho_c$.
Clock covariance and Holevo data processing give $a(\sigma)\leq a(\rho)$;
pinching preserves $X$ and $Q_\gamma$.
For an actual word $C$ the identity and generator weights give
\[
 \chi(C,\sigma):=\|C\sqrt\sigma C^*-\sqrt\sigma\|_F
       \leq |C|c_g\sqrt a.
\]
Tracing to the clock bounds each cyclic root energy by $c_g^2a$.
The torus Poincar\'e inequality and differences of squares imply
$\operatorname{TV}(p,\mathrm{unif})\leq c_gM\sqrt a$.
For $M\geq4096$ put $x=\lfloor M/4\rfloor$, $R=x+8$, and
\[
 G=\{(i,j):0\leq i\leq M-x-12,\quad
               1\leq j\leq M-3,\quad i+j\geq M-x+16\}.
\]
Its size is $(M-x-11)(M+x-48)/2$, so its density is at least
$15/32-40/M\geq0.45$. Restricting the original anchor charge gives
\begin{equation}\label{eq:houghton-wedge-charge}
 q_G:=\sum_{c\in G}p_c\Tr\rho_cQ_\gamma
       \geq1/30-c_gM\sqrt a.
\end{equation}
Thus $M^2a\leq\varepsilon_0:=1/(120c_g)^2$ implies $q_G\geq1/40$.

We need a stronger version of the preceding interior-chain geometry.
At $c=(i,j)\in G$, horizontal prefixes of length at most $R-1$
have edge starts at most $M-6$ and endpoints at most $M-5$.
They meet neither the seam nor the preliminary exact $A$ move.
The first reverse cycle has length $r_0\geq x+17=R+9$.
In inverse traversal the tracked indices and cycle lengths increase
together, so there is no wrap, even though later cycles can be
shorter than $R+4$; the last cycle has length at least eleven.
The row bound $j\leq M-3$ excludes the exceptional $T$-source row.
As before, bulk and $T$ eligibility require both incident plaquettes,
and are respectively $i+j\leq t+M-8$ and $i+j\leq M-3$.
Pulled-back planes are disjoint fixed pairs $(s,s+M+i)$.
A bulk pair can rotate at prefix step $q$ only for $q\leq s-7$,
and a $T$ pair only for $q\leq s-6$. Its rotation is therefore
stable before a reflection normal uses $e_s$.
One common product $D_c$ of these commuting rotations gives
orthonormal points $f_s=D_c^*e_s$ and reflection normals
$(f_{q+1}-f_{q+2})/\sqrt2$, for $0\leq q<R$.
Thus the chain is exactly $S_{R+1}$, with
$f_1,f_2,f_3=e_1,e_2,e_3$.
All the same pairs support the field $w_c$, so $D_c$ commutes with it.

Choose $r\in\{0,\ldots,7\}$, set $s=x+r$ and
$t=x+r-M+1+i+j$. Then $17\leq t\leq M-7$ and this bulk
region is active, with pair $(s,s+M+i)$.
Its untruncated triangle count gives
\[
 A_t=(M-1)^2-h(h+1)/2,\qquad h=M+3-t.
\]
Triangular residues modulo four repeat $0,1,3,2,2,3,1,0$ over
eight consecutive values. Since $(M-1)^2$ is zero or one modulo
four, one choice has $4\mid A_t$.
Put $p=A_t/4\leq M^2/4$, relabel $f_s$ as point four, and take
$f_5=D_c^*e_{s+M+i}$. On these five points $W_c^p$ is
$g_4=\exp(\pm i\pi P_{f_4-f_5}/4)$, where $P$ projects onto
the normalized difference. All other phases are outside the five
points and commute with $S_4$ on the first four.
The swap of low point $i_0\leq3$ with $s$ is the actual word
\[
 C_{i_0,s}=\sfa^{i_0-1}(\alpha\sfa)^{s-i_0-1}\alpha
       (\sfa^{-1}\alpha)^{s-i_0-1}\sfa^{-(i_0-1)}.
\]
Adjacent products collapse to this word of length at most
$4s\leq M+28\leq2M$. With $C_4=I$, there are only $24$
nontrivial global choices. Hence conditional selection costs
\begin{equation}\label{eq:houghton-word-congestion}
 \sum_{c\in G}p_c\sum_{i=1}^3\chi(C_i^{r(c)},\rho_c)^2
       \leq96c_g^2M^2a.
\end{equation}

Here is the fixed-group gap input. In the canonical five-point basis
put $v_i=(e_i-e_5)/\sqrt2$, $g_i=e^{i\pi P_{v_i}/4}$,
and take the twelve matrices $U_{ij}=g_ig_j^*$ and
$g_iU_{ij}g_i^*$, $i<j$. They fix the all-ones vector and lie
in $SU(H_0)$, $H_0=\{z:\sum z_i=0\}$, of dimension four.
On $E_{ij}=\operatorname{span}(v_i,v_j)$ the nontrivial
eigenvalues of $U_{ij}$ are $e^{\pm i\phi}$, with
$2\cos\phi=(2+3\sqrt2)/4$. Its minimal polynomial is
$z^2-z-7/8$, so it is not an algebraic integer.
Thus $\phi/\pi$ is irrational and powers give a continuous
$SU(2)$ circle on $E_{ij}$. Its $g_i$-conjugate is nonparallel:
$[g_i,U_{ij}]=g_i[g_i,g_j^*]\neq0$; an antiparallel axis is
impossible since $\operatorname{Ad}g_i$ has no eigenvalue $-1$.
The two circles generate $SU(2)(E_{ij})$.
The six real skew matrices $v_iv_j^T-v_jv_i^T$ span
$\mathfrak{so}(4)$. The same plane algebras contain the nonzero
imaginary symmetric traceless $i(P_{v_1}-P_{v_2})$.
Orthogonal conjugation diagonalizes it to a nonzero multiple of
$i\operatorname{diag}(1,-1,0,0)$; permutations and plane rotations
span all nine imaginary symmetric traceless directions.
Together these span $\mathfrak{su}(4)$, proving density.
An algebraic orthonormal basis of $H_0$ makes all entries algebraic.
Bourgain--Gamburd's theorem~\cite[Theorem~1]{BourgainGamburd2012}
and Peter--Weyl give a fixed $\delta>0$ such that, in every
continuous representation,
\[
 \|(I-P_{\rm inv})\xi\|^2
       \leq\delta^{-1}\sum_h\|(h-I)\xi\|^2.
\]
We use the unnormalized symmetric gap. The negative-angle list is
the complex conjugate and has the same gap.

For the actual gates $G_i=C_iW_c^pC_i^*$ the common outside
factor cancels in $G_iG_j^*$ and $G_i(G_iG_j^*)G_i^*$.
Quantization gives continuous representations of this same $SU(4)$
for every $k$, with arbitrary spectators.
The original three-cycle $\gamma$ belongs to it, so $Q_\gamma$
kills its invariant vectors. Apply the gap to left multiplication
on Hilbert--Schmidt matrices, not to conjugation.
For $\xi=\sqrt{\rho_c}$, $E=\|(W_c-I)\xi\|_F$ and
$\chi_i=\|C_i\xi C_i^*-\xi\|_F$, moving $C_i^*$ across
$\xi$ gives $\|(G_i-I)\xi\|_F\leq pE+2\chi_i$.
Products telescope in left-action norm and inverses have the same
energy. The two listed gates have respective bounds
$2pE+2\chi_i+2\chi_j$ and $4pE+6\chi_i+2\chi_j$.
Three-term Cauchy and summation give
\[
 \sum_h\|(h-I)\xi\|_F^2
       \leq360p^2E^2+360\sum_{i=1}^3\chi_i^2.
\]
Since the selected power obeys $p(c)^2\leq M^4/16$ pointwise,
\eqref{eq:houghton-word-congestion} yields
\[
 q_G\leq\delta^{-1}\left[(45/2)M^4\Tr\rho X
                               +34560c_g^2M^2a\right].
\]
With $A=34560c_g^2$, $B=(45/2)(225d/488)$,
\eqref{eq:houghton-wedge-charge} proves the bound at $M\geq4096$
with $c_{\rm large}=\min\{\varepsilon_0,\delta/[40\max(A,B)]\}$.
For $48\leq M<4096$, Theorem~\ref{thm:fixed-clock} gives
$M^2a+M^4\ell\geq M^2c_M$. The minimum of these finitely
many positive constants and $c_{\rm large}$ is $c_0>0$.
No numerical gap constant is asserted.

Finally put $t_0=\min\{n^{1/3},\ell^{-1/4}\}$ and
$\beta=\min\{1,(c_0/2)^{1/4}\}$.
If $M\geq\beta t_0$, Theorem~\ref{thm:houghton-entropy} gives
$S+na+1\geq c_{\rm ent}\beta t_0$.
Otherwise $M^4\ell<c_0/2$, so
$na\geq c_0n/(2M^2)>[c_0/(2\beta^2)]n/t_0^2
\geq[c_0/(2\beta^2)]t_0$, since $t_0^3\leq n$.
This proves the second inequality. It gives the stated $n^{1/4}$
floor at $\ell\leq K/n$, and a cube-root lower at
$\ell\leq n^{-4/3}$, without asserting a matching zero-leakage upper.
\end{proof}

\subsection{Zero production and the actual stationary algebra}
\label{app:houghton-stationary}

Let $\mathcal E_M$ be Haar averaging over the closed group generated
by the \emph{actual} slot unitaries on the counted bath. Its range is
their commutant. It commutes with clock pinching $\Delta_c$, since
each slot has a fixed clock translation and a blockwise internal action.
This averaging is an analytical conditional expectation.

\begin{corollary}[Stationary sources at arbitrary flavour count]
\label{cor:houghton-stationary}
There are constants $b,\delta_*,c>0$, independent of $M\geq48$ and
$k\geq1$, with the following properties. Put $h=M^{-3}$,
$F=\min(X,hI)$ and $q_\eta=\Tr\eta Q_\gamma$.
Every actual stationary state $\eta=\mathcal E_M\eta$ obeys
\begin{equation}\label{eq:stationary-capped-charge}
 \Tr\eta F/h\geq b(q_\eta-11/20).
\end{equation}
If a physical source has $u\leq1/16$ and actual $a=0$, then
\[
 \ell\geq cM^{-3},\qquad S(\rho)+1\geq cM,
 \qquad S(\rho)+1\geq c\ell^{-1/3}.
\]
More generally, $u\leq1/16$ and
\[
 \tfrac12\|\Delta_c\rho-\mathcal E_M(\Delta_c\rho)\|_1
       \leq\delta_*
\]
imply $\ell\geq cM^{-3}$. If also $\ell\leq K/n$ for fixed
$K>0$, then $S(\rho)+na(\rho)+1\geq c_Kn^{1/3}$.
Entropy and production always refer to the original source.
\end{corollary}

\begin{proof}
The Holevo identity and the identity slot show that $a=0$ holds
exactly when every positive-weight slot fixes $\rho$ by conjugation.
Thus $\rho$ commutes with every actual word. Clock pinching retains
stationarity, $X$ and $Q_\gamma$. The two clock translations force
each pinched clock mass to be $1/M^2$. For any stationary $\eta$,
write $\Delta_c\eta=M^{-2}\bigoplus_c\eta_c$.

For $M\geq4096$, use the exact stabilized logical-point geometry
above, restricted to the smaller wedge of
Appendix~\ref{app:houghton-moderate}. Put $x=\lfloor M/4\rfloor$,
$z=\lfloor M/256\rfloor$ and
\[
 G=\{0\leq i\leq M-x-12,\ 1\leq j\leq M-3,
                      \ i+j\geq M-x+16+2z\}.
\]
Its density is at least $0.45$. Choose even
$b_0\in[x-z,x-z+1]$, $J=\lfloor(z-1)/8\rfloor$ and
$b_j=b_0+8j$, $0\leq j<J$. Then $M/8192\leq J\leq M/2048$.
In each eight-index interval choose $s_j$ so that the actual region
$t=s_j-M+1+i+j$ has $4\mid A_t$.
The triangular residues used above guarantee this choice.
These active regions have $M^2/3\leq A_t\leq M^2$.
Their pairs $(f_{s_j},D_c^*e_{s_j+M+i})$ are disjoint from one
another and from the common anchor modes $f_1,f_2,f_3$.
The usual-representation angles have absolute value $\pi/A_t$;
taking adjoints reverses the written right-action sign.

The rational power $A_t/4$ and the three low-to-selected swaps
give the fixed local $SU(4)$ of the preceding proof on the sum-zero
space $H_j$ of those five modes. Outside phases cancel in its
generators. Stationarity gives exact commutation with these groups;
no word-length or congestion estimate is needed.
They generate $SU(H)$ on the sum-zero space of all $3+2J$ modes,
of dimension $2J+2$. To verify this, choose a unit vector $e$ in the
common two-dimensional anchor sum-zero plane. Each local Lie algebra
contains both off-diagonal skew-Hermitian generators between $e$ and
every vector in $H_j\cap e^\perp$. The spaces $H_j$ span $H$:
the two common anchor differences and the two independent differences
from each selected pair to an anchor give a spanning set.
Thus these local perpendicular spaces span $H\cap e^\perp$.
Linearity and commutators give every off-diagonal direction and every
traceless diagonal direction of $\mathfrak{su}(H)$.

On $H$, the field has $J$ orthogonal active difference modes $v_j$
with angles $\theta_j$, and fixes a complementary $(J+2)$-space.
Pair the $v_j$ with $J$ orthogonal fixed modes. The resulting product
$SU(2)$ groups lie in $SU(H)$. An invariant conditional state has
the same nontrivial spin mass $p_c$ on every such plane, by conjugacy.
The anchor three-cycle lies in an anchor $SU(2)$, so
$q_c\leq p_c$. Arbitrary flavour and multiplicity correlations are
allowed in its flat product-spin decomposition.

Apply the capped spin estimate of Appendix~\ref{app:houghton-stability}
to the commutator of the full field with the product of the paired
quarter turns $v_j\mapsto w_j$, $w_j\mapsto-v_j$. Each quarter turn lies
in $SU(2)\subset SU(H)$, unlike a swap, whose determinant on $H$ is
$-1$, and it produces the same commutator. Outside phases cancel. Its
transfer to $F$ is valid because the stationary conditional state
commutes with both the full field and the quarter turns, making the
corresponding functional tracial.
Here $\sum_j\theta_j^2\leq9\pi^2/(2048M^3)<h$ and
$|\theta_j|\leq\pi/4$. Consequently
\[
 \Tr\eta_c F\geq\frac{p_c}{2048}\sum_j\theta_j^2
      \geq\frac{q_cJ\pi^2}{2048M^4}.
\]
Uniform clock mass and positivity allow discarding clocks outside $G$.
Their lost charge is at most $11/20$, giving
\eqref{eq:stationary-capped-charge} with
$b_{\rm large}=\pi^2/(2048\cdot8192)$.
For $48\leq M<4096$, $X\leq4I$ gives $F\geq hX/4$.
Equation~\eqref{eq:fixed-clock-operator} at stationarity gives
$q_\eta\leq C_M\Tr\eta X$. Choose
\[
 b=\min\left\{b_{\rm large},\min_{48\leq M<4096}(4C_M)^{-1}\right\}>0.
\]
This is an existence constant; the finite-clock gaps are not evaluated.

Put $\beta=b/80$ and $\delta_* =\min\{1/48,\beta/2\}$.
For $\sigma=\Delta_c\rho$ and $\eta=\mathcal E_M\sigma$,
the stated trace distance transfers only the anchor charge:
$q_\eta\geq7/12-1/48=9/16$.
Hence $\Tr\eta F\geq\beta h$ and, since $0\leq F\leq hI$,
\[
 \Tr\rho X=\Tr\sigma X\geq\Tr\sigma F
       \geq(\beta-\delta_*)h\geq\beta h/2.
\]
Equation~\eqref{eq:actual-houghton-leakage} proves the leakage lower.
At $a=0$ the entropy theorem gives $S+1\geq c_{\rm ent}M$;
combining with $\ell\geq cM^{-3}$ gives its third inequality.
For $\ell\leq K/n$, the leakage bound forces
$M\geq(cn/K)^{1/3}$, and
Theorem~\ref{thm:houghton-entropy} gives the final cost bound.
No entropy of $\sigma$ or $\eta$ is substituted for $S(\rho)$.
\end{proof}

The condition concerns the actual stationary algebra, which need not
coincide with the collective symmetry algebra. Small entropy production
does not automatically supply its required trace proximity uniformly
as $M$ grows. The corollary is about this literal family and does not
give an unrestricted-source or all-realization lower bound.

\subsection{The unresolved sharp quantitative estimate}

Write $\mathcal D=\sum_sw_s(\operatorname{Ad}V_s-I)^*
(\operatorname{Ad}V_s-I)$,
$\mathcal X=(L_X+R_X)/2$ and
$\mathcal Q=(L_{Q_\gamma}+R_{Q_\gamma})/2$ on Hilbert--Schmidt space.
The concrete sufficient estimate
\begin{equation}\label{eq:open-killed-laplacian}
 M^2\mathcal D+M^3\,\operatorname{polylog}(M)\mathcal X
       \ \geq\ \kappa\mathcal Q
\end{equation}
uniformly in $k$, with $\kappa>0$ independent of $M$, is open.
Theorem~\ref{thm:houghton-entropy} supplies the entropy half for all
coherent sources in this literal geometry at growing $M$.
Theorem~\ref{thm:houghton-flux} gives a weaker physical-source
coercivity with $M^4\ell$, not this sharper operator inequality.
Its fixed rational power $p=O(M^2)$ explains the extra factor of $M$.
Theorem~\ref{thm:houghton-moderate} removes that factor for physical
sources with $k\leq c_f\sqrt M$, using many small-angle blocks.
It does not prove the displayed operator inequality uniformly in $k$.
Unrestricted flavour growth, all-round scope and optimized universal
comparison remain unproved.

There is a useful exact reduction for testing the operator estimate.
For $c_D,c_X\geq0$, positivity of
$\mathcal H=c_D\mathcal D+c_X\mathcal X-\lambda\mathcal Q$
on all complex Hilbert--Schmidt matrices is equivalent to positivity
on positive matrices. Indeed $\mathcal H$ commutes with adjoint,
so a negative witness splits into a negative Hermitian witness.
For Hermitian $A$, replacing $A$ by $|A|$ keeps $A^2$ and hence
both potential energies. It lowers the Dirichlet energy because
\[
 \||A|-|B|\|_F\leq\|A-B\|_F\quad(A=A^*,\ B=B^*).
\]
To check this inequality, expand in eigenbases: each spectral
overlap term has $|a||b|\geq ab$. Apply it with
$B=V_sAV_s^*$. Thus a negative positive witness exists.

The operator commutes with clock pinching and with each flavour's
degree pinching. A positive ground witness survives pinching
(its trace is nonzero) and stays in the ground eigenspace.
A nonzero degree-tuple block is therefore a clock-diagonal
positive witness. This is a spectral-test reduction, not an
entropy-preserving source twirl: pinching can increase source
entropy, and the operator witness need not satisfy all the
channel's coarse-correctness constraints.
